% !TeX root = ../thesis.tex
\chapter{THIẾT KẾ VÀ THỰC HIỆN PHẦN MỀM}

\section{Kiến trúc phần mềm ROS 2}

\subsection{Tổng quan kiến trúc node}

Hệ thống phần mềm được triển khai trên nền tảng ROS~2 Humble \cite{ros2docs}, sử dụng kiến trúc xử lý phân tán gồm hai node tính toán được phân vùng theo chức năng:

\begin{itemize}
    \item \textbf{Gait Generator Node (Node Điều khiển dáng đi):} Quản lý toàn bộ quá trình tính toán quỹ đạo tuần hoàn, giải bài toán động học nghịch và gửi lệnh điều khiển đến các cơ cấu chấp hành qua bus nối tiếp. Node này được tổ chức thành các phân hệ chức năng: Phân hệ tạo dáng đi (quản lý trạng thái), Phân hệ động học nghịch (IK solver), Phân hệ nội suy quỹ đạo (đa thức bậc 5) và Phân hệ giao tiếp phần cứng (UART).
    
    \item \textbf{Balance Controller Node (Node Cân bằng tư thế):} Xử lý bất đồng bộ việc thu thập dữ liệu quán tính tần số cao, xử lý tín hiệu đa giai đoạn và tính toán bù PID. Node này hoạt động như một tiến trình độc lập để đảm bảo chu kỳ điều khiển không bị gián đoạn.
\end{itemize}

\subsection{Kiến trúc phân tầng và phương pháp thiết kế Sim-to-Real}

Điểm nổi bật trong kiến trúc phần mềm ROS~2 của hệ thống là khả năng chuyển đổi liền mạch giữa môi trường mô phỏng (Gazebo) và phần cứng thực tế (Hardware) mà không yêu cầu chỉnh sửa thuật toán điều khiển cốt lõi (tư tưởng Sim-to-Real). Để đạt được điều này, hệ thống được thiết kế theo cấu trúc phân tầng (layered architecture), tách biệt rõ ràng các khối xử lý thuật toán khỏi các giao diện điều khiển thiết bị vật lý.

Trong môi trường mô phỏng (Hình~\ref{fig:ros2_arch_sim}), kiến trúc phần mềm được chia thành hai phân vùng chính:
\begin{itemize}
    \item \textbf{Tầng điều khiển phần mềm (ROS 2 Software Control):} Chứa hai node thuật toán cốt lõi là \texttt{Balance Controller Node} và \texttt{Gait Generator Node}. Hai node này giao tiếp nội bộ với nhau qua các topic \texttt{/posture/correction} để chia sẻ tín hiệu bù thăng bằng.
    \item \textbf{Tầng mô phỏng (Gazebo Simulation Environment):} Hoạt động như một robot ảo, trong đó \texttt{imu\_plugin} đóng vai trò xuất bản dữ liệu quaternion thay cho cảm biến thật, còn \texttt{controller\_manager} đóng vai trò quản lý mô-men và góc khớp ảo dựa trên lệnh \texttt{/effort\_controller} từ bộ tạo dáng đi.
\end{itemize}

\begin{figure}[H]
    \centering
    \includegraphics[width=0.95\textwidth]{architecture/simulation_rqt_graph.png}
    \caption{Sơ đồ kiến trúc khối ROS~2 của hệ thống trong môi trường mô phỏng (Gazebo)}
    \label{fig:ros2_arch_sim}
\end{figure}

Ngược lại, khi chuyển sang vận hành trên robot thực tế (Hình~\ref{fig:ros2_arch_hw}), hệ thống thay thế toàn bộ Tầng mô phỏng bằng hai tầng vật lý chuyên dụng, trong khi vẫn \textbf{giữ nguyên} nguyên bản Tầng thuật toán cốt lõi (Tầng 1). Kiến trúc thực tế bao gồm ba tầng rõ rệt:
\begin{itemize}
    \item \textbf{Tầng 1 -- ROS 2 Core Algorithms:} Chứa các node giải quyết bài toán động học, nội suy quỹ đạo và điều khiển PID (tái sử dụng 100\% mã nguồn từ mô phỏng).
    \item \textbf{Tầng 2 -- Hardware Interface Drivers:} Đóng vai trò cầu nối, bao gồm \texttt{imu\_sensor\_node} để chuyển đổi luồng I2C thành message ROS~2 tiêu chuẩn, cùng với \texttt{serial\_driver\_a} và \texttt{serial\_driver\_b} làm nhiệm vụ dịch các mảng góc khớp mục tiêu thành gói tin bus servo thông qua thư viện SCServo SDK.
    \item \textbf{Tầng 3 -- Physical Hardware:} Gồm các linh kiện điện tử vật lý là cảm biến BNO055 IMU và chuỗi 12 động cơ ST3215 Servo giao tiếp theo chuẩn UART/RS485.
\end{itemize}

Sự phân tách ba tầng này giúp hệ thống đạt độ tin cậy vận hành rất cao, đồng thời cho phép việc bảo trì hay nâng cấp cảm biến, động cơ trong tương lai chỉ gói gọn ở việc thay đổi driver ở Tầng 2 mà không ảnh hưởng tới chất lượng thuật toán.

\begin{figure}[H]
    \centering
    \includegraphics[width=0.95\textwidth]{architecture/hardware_rqt_graph.png}
    \caption{Sơ đồ kiến trúc ROS~2 phân tầng của hệ thống trên phần cứng thực tế}
    \label{fig:ros2_arch_hw}
\end{figure}

\subsection{Các trạng thái vận hành}

Hệ thống quản lý hai trạng thái vận hành rời rạc thông qua cổng logic boolean:

\begin{itemize}
    \item \textbf{HOME:} Cấu hình tĩnh không có vector tịnh tiến tuần hoàn. Bộ điều khiển tính toán các điều chỉnh đồng nhất trong không gian khớp, ánh xạ cứng đến tư thế nghỉ. Sai lệch góc cơ sở được chuyển đổi tỷ lệ thành các bù mô-men đối kháng nhằm hiệu chỉnh tư thế liên tục. Một khoảng trễ khởi động $T_{delay}$ giây có thể cấu hình giúp ngăn hiệu chỉnh sớm trong giai đoạn quá độ khởi tạo hệ thống.
    
    \item \textbf{GAIT:} Trạng thái động kích hoạt khi có lệnh tịnh tiến mặt phẳng. Robot thực hiện quỹ đạo đa thức tuần hoàn sử dụng các mảng dịch pha khả cấu hình để tạo dáng đi trot (với LF, RB: 0 và LB, RF: 0,5), walk (với LF: 0, RB: 0,25, RF: 0,5, LB: 0,75) và các dáng đi khác. Các hiệu chỉnh không gian từ phản hồi quán tính chỉ được áp dụng cho chân ở pha chống, trong khi chân ở pha bay duy trì bù bằng không để đảm bảo khoảng trống sải bước.
\end{itemize}

\subsection{Cấu trúc các ROS 2 packages}

Toàn bộ hệ thống phần mềm được tổ chức thành 6 packages ROS~2:

\begin{table}[H]
\centering
\caption{Cấu trúc các ROS 2 packages trong hệ thống}
\label{tab:ros2_packages}
\renewcommand{\arraystretch}{1.3}
\begin{tabular}{|p{0.25\linewidth}|p{0.12\linewidth}|p{0.49\linewidth}|}
\hline
\textbf{Package} & \textbf{Ngôn ngữ} & \textbf{Chức năng} \\ \hline
\texttt{leg\_controller} & Python & Tạo dáng đi, IK solver, quy hoạch quỹ đạo, giao tiếp servo \\ \hline
\texttt{balance\_controller} & Python & Thu thập IMU, xử lý tín hiệu, PID cân bằng, chẩn đoán \\ \hline
\texttt{serial\_driver} & C++ & Driver giao tiếp với bus servo ST3215 (thư viện SCServo) \\ \hline
\texttt{simulation} & Python & Mô hình URDF, cấu hình Gazebo, RViz, launch mô phỏng \\ \hline
\texttt{launch\_main} & CMake & Launch files cho hệ thống thực \\ \hline
\end{tabular}
\end{table}

\subsection{Chi tiết các module phần mềm}

\textbf{Phân hệ điều khiển trung tâm (Gait Generator)} là phân hệ cốt lõi của hệ thống, quản lý máy trạng thái (state machine) điều khiển toàn bộ vòng đời vận hành của robot. Phân hệ này thực hiện các chức năng:

\begin{itemize}
    \item \textbf{Quản lý trạng thái:} Robot chuyển đổi giữa các trạng thái: INIT (khởi tạo) $\rightarrow$ ZERO (khởi tạo tư thế đứng) $\rightarrow$ TROT/WALK/PUSHUP/SWAY (dáng đi) $\rightarrow$ ROBOTOFF (tắt an toàn). Mỗi trạng thái có chu trình cập nhật độc lập để xử lý logic trong một khung thời gian điều khiển (frame).
    \item \textbf{Tạo quỹ đạo:} Gọi thuật toán nội suy quỹ đạo để tạo quỹ đạo hình chữ D cho mỗi chân, sau đó giải phương trình động học nghịch (IK) để chuyển đổi từ tọa độ Đề-các sang không gian góc khớp.
    \item \textbf{Áp dụng hiệu chỉnh IMU:} Nhận tín hiệu bù từ Balance Controller qua hai topic \texttt{/posture/}\allowbreak\texttt{home\_correction} và \texttt{/posture/}\allowbreak\texttt{gait\_correction}, áp dụng vào góc khớp (HOME) hoặc tọa độ bàn chân (GAIT).
    \item \textbf{Xuất lệnh servo:} Gửi góc điều khiển đến phân hệ giao tiếp phần cứng để chuyển đổi từ góc khớp phân tích sang góc vật lý của servo và truyền qua giao thức UART.
\end{itemize}

\textbf{Phân hệ giao tiếp phần cứng (Serial Publish)} đảm nhận việc chuyển đổi giữa hai chế độ vận hành (mô phỏng và phần cứng thực) thông qua cấu hình môi trường:

\begin{itemize}
    \item \textbf{Chế độ mô phỏng:} Xuất bản lệnh mô-men (effort) lên topic \texttt{/effort\_controller/commands} để điều khiển robot trong Gazebo thông qua plugin \texttt{gazebo\_ros2\_control}.
    \item \textbf{Chế độ phần cứng:} Chuyển đổi góc từ động học nghịch sang độ phân giải của servo bằng các hàm ánh xạ riêng biệt cho từng chân (xử lý ngược chiều quay do lắp đặt đối xứng), sau đó đóng gói dữ liệu và xuất bản lên các topic giao tiếp serial.
\end{itemize}

\subsection{Bảng giao tiếp topic ROS~2}

Bảng~\ref{tab:ros2_topics} liệt kê các topic chính trong hệ thống, phân loại theo chức năng.

\begin{table}[H]
\centering
\caption{Các topic ROS~2 chính trong hệ thống}
\label{tab:ros2_topics}
\renewcommand{\arraystretch}{1.3}
\begin{tabular}{|p{0.32\linewidth}|p{0.12\linewidth}|p{0.42\linewidth}|}
\hline
\textbf{Topic} & \textbf{Kiểu dữ liệu} & \textbf{Mô tả} \\ \hline
\texttt{/imu/data} & Imu & Dữ liệu quaternion từ BNO055 \\ \hline
\texttt{/balance/enable} & Bool & Bật/tắt bộ PID cân bằng \\ \hline
\texttt{/balance/mode} & String & Chuyển chế độ HOME/GAIT \\ \hline
\texttt{/balance/gait\_frame} & Int32 & Chỉ số frame dáng đi hiện tại \\ \hline
\texttt{/posture/home\_correction} & Vector3 & Tín hiệu bù PID chế độ HOME \\ \hline
\texttt{/posture/gait\_correction} & Vector3 & Tín hiệu bù PID chế độ GAIT \\ \hline
\texttt{/posture/measurement} & Vector3 & Góc roll/pitch đã lọc \\ \hline
\texttt{/servo\_commands\_a} & Float64[] & Lệnh góc cho Driver A (6 servo) \\ \hline
\texttt{/servo\_commands\_b} & Float64[] & Lệnh góc cho Driver B (6 servo) \\ \hline
\end{tabular}
\end{table}

\subsection{Cơ chế chẩn đoán thời gian thực}

Hệ thống tích hợp cơ chế chẩn đoán (diagnostics) cho phép giám sát trực tiếp các biến nội bộ của bộ điều khiển PID trong thời gian thực thông qua công cụ PlotJuggler. Mỗi giai đoạn trong pipeline PID (đo lường $\rightarrow$ sai lệch $\rightarrow$ P/I/D $\rightarrow$ bão hòa $\rightarrow$ lọc) đều có topic chẩn đoán riêng, cho phép:

\begin{itemize}
    \item Quan sát trực tiếp từng thành phần P, I, D và tổng PID
    \item Phát hiện hiện tượng windup hoặc bão hòa đầu ra
    \item So sánh tín hiệu đo (measurement) với điểm đặt (setpoint) để đánh giá chất lượng bám
    \item Ghi lại dữ liệu dưới dạng file CSV để phân tích offline
\end{itemize}

Tổng cộng, hệ thống xuất bản 22 topic chẩn đoán cho chế độ HOME (9 tín hiệu $\times$ 2 trục + 4 tín hiệu áp dụng) và 18 topic cho chế độ GAIT (9 tín hiệu $\times$ 2 trục), cung cấp khả năng quan sát toàn diện luồng xử lý điều khiển.

\section{Quy trình Homing và Calibrate}

\subsection{Calibrate servo}

Trước khi vận hành, mỗi servo cần được calibrate để đảm bảo góc $0^\circ$ cơ học trùng với góc $0^\circ$ trong phần mềm. Quy trình calibrate bao gồm:

\begin{enumerate}
    \item Đặt tất cả chân robot ở tư thế tham chiếu (chân duỗi thẳng, vuông góc với thân).
    \item Đọc vị trí feedback của từng servo qua thư viện SCServo SDK.
    \item Ghi giá trị offset ($\Delta_{calib}$) --- sai lệch giữa vị trí thực tế và vị trí lý thuyết tại góc $0^\circ$.
    \item Lưu offset vào tập tin cấu hình của hệ thống, áp dụng tự động trong các lần khởi chạy tiếp theo.
\end{enumerate}

\subsection{Quy trình Homing}

Quy trình khởi tạo tư thế đứng (homing) được thực hiện tự động mỗi khi robot khởi động:

\begin{enumerate}
    \item \textbf{Đọc vị trí hiện tại:} Lấy giá trị góc feedback từ tất cả 12 servo ($\theta_{current,i}$).
    \item \textbf{Tính tư thế đích:} Giải IK cho tư thế đứng chuẩn, tạo ra bộ góc đích $\theta_{target,i} = (0^\circ, 0^\circ, 90^\circ)$ cho mỗi chân.
    \item \textbf{Nội suy tuyến tính:} Trong $N_{homing} = 200$ frame, mỗi góc được nội suy tuyến tính:
    \begin{equation}
        \theta_i(k) = \theta_{current,i} + \frac{k}{N_{homing}} (\theta_{target,i} - \theta_{current,i}), \quad k = 0, 1, \ldots, N_{homing}
    \end{equation}
    \item \textbf{Gửi lệnh:} Tại mỗi frame, gửi bộ 12 góc đã nội suy đến servo qua UART.
\end{enumerate}

Phương pháp nội suy tuyến tính đảm bảo chuyển động mượt mà từ bất kỳ vị trí khởi đầu nào, tránh hiện tượng giật hoặc va đập. Tổng thời gian khởi tạo: $200 \times 3 = 600$ ms $= 0{,}6$ giây.

